\documentclass[11pt,a4paper,twoside]{book}
\usepackage[utf8]{inputenc}
% \usepackage[T1]{fontenc}
\usepackage{graphicx}
\usepackage{tabularx}
\usepackage{array}
\usepackage{twoopt}
\usepackage{enumitem}
\usepackage{calc}
\usepackage[hyphens]{url}
\usepackage{longtable}
\usepackage{hyperref}
\usepackage{color}
\usepackage{ifthen}
\usepackage{xifthen}
\hypersetup{colorlinks=true,linkcolor=blue,citecolor=blue,pdfstartview={XYZ null null 1.00}}
\usepackage{amsmath}



\def\ifempty#1{%
\edef\ContentsOfList{\the#1}%
\edef\CompareTo{}%
\ifx\ContentsOfList\CompareTo}



%%%%%%%%%%%%%%%%%%%%5
\makeatletter
\def\nobreakhline{%
  \noalign{\ifnum0=`}\fi
    \penalty\@M
    \futurelet\@let@token\LT@@nobreakhline}
\def\LT@@nobreakhline{%
  \ifx\@let@token\hline
    \global\let\@gtempa\@gobble
    \gdef\LT@sep{\penalty\@M\vskip\doublerulesep}% <-- change here
  \else
    \global\let\@gtempa\@empty
    \gdef\LT@sep{\penalty\@M\vskip-\arrayrulewidth}% <-- change here
  \fi
  \ifnum0=`{\fi}%
  \multispan\LT@cols
     \unskip\leaders\hrule\@height\arrayrulewidth\hfill\cr
  \noalign{\LT@sep}%
  \multispan\LT@cols
     \unskip\leaders\hrule\@height\arrayrulewidth\hfill\cr
  \noalign{\penalty\@M}%
  \@gtempa}
\makeatother
%%%%%%%%%%%%%%%%%%%%%%%%%%%%%%%%%5

% Environment used for option descriptions
% \newenvironment{optiontable}{\noindent\begin{small}\begin{tabular*}{\textwidth}{p{0.4\textwidth}llp{0.4\textwidth}}name & type & default & description \\\hline}{\end{tabular*}\end{small}}
% \newenvironment{optiontable}{\setlength\LTleft{0pt}\setlength\LTright{0pt}\noindent\begin{small}\begin{longtable}{p{0.3\textwidth}llp{0.4\textwidth}}name & type & default & description \\\hline}{\end{longtable}\end{small}}
\newenvironment{optiontable}{\setlength\LTleft{0pt}\setlength\LTright{0pt}\noindent\begin{small}\begin{longtable}{p{0.4\textwidth}p{0.5\textwidth}}name & description \\*\nobreakhline}{\hline\end{longtable}\end{small}}
% \newcommand{\option}[4]{\path{--#1} & #2 & #3 & #4 \\}
% \newcommand{\option}[4]{\path{--#1} & #3 & #4 \\}
% \newcommand{\option}[4]{\ifempty{#1} empty\else nonempty\fi}
\newcommand{\option}[4]{\path{#1}\ifthenelse{\isempty{#3}}%
    {}% if #1 is empty
    {\mbox{(=#3)}}% if #1 is not empty
    & #4 \\}


\newcommand{\optiontitle}[1]{\subsubsection*{#1}\vspace*{-1em}}

\newcommand{\phaistosversion}{@MANUAL_VERSION@}
\newcommand{\phaistosauthors}{@MANUAL_AUTHORS@}
\newcommand{\phaistosbibtexfiles}{@phaistos_BIB_FILES_FORMATTED@}

\begin{document}

\newcommand{\notification}[1]{
\noindent
\begin{minipage}{\textwidth}
  \definecolor{lightgray}{gray}{0.85} %
  \definecolor{darkgray}{gray}{0.25} %
  \centering %
  \vspace*{.5em} %
  \setlength\fboxrule{1.5pt} %
  \fcolorbox{darkgray}{lightgray}{\parbox{.9\textwidth}{#1}} %
  \setlength\fboxsep{0pt} %
  \vspace*{.5em} %
\end{minipage} %
}



\begin{titlepage}
  
  \begin{center}

    \vspace*{5em}
    
    \begin{Huge}
      Phaistos User Manual
    \end{Huge}
    
    \vspace*{1em}
    
    \begin{LARGE}
      Version \phaistosversion
    \end{LARGE}

    \vspace*{6em}

    \includegraphics[width=.5\textwidth]{phaistos_logo_noname.pdf}
    
    \vspace*{6em}
    
    \begin{Large}
      \phaistosauthors
    \end{Large}
    
  \end{center}
  
\end{titlepage}

\newpage

\tableofcontents

\newpage

\chapter{Introduction}
\label{cha:introduction}
Phaistos is a molecular modelling software package for 
simulation of proteins. The distinctive features of Phaistos 
are the use of the Markov chain Monte Carlo (MCMC) approach to 
molecular simulation, and the use of sophisticated probabilistic models
of protein structure in continuous space. Phaistos 
can be used for molecular simulation using
empirical energy functions such as OPLS-AA or Profasi, or 
for probabilistic, Bayesian inference of protein structure based on
a likelihood and conformational priors. The latter applications 
include for example the inference of protein structure using a
likelihood that brings in NMR or SAXS data, and priors that 
takes into account the general features of protein structure. 
Phaistos covers a wide spectrum of applications, as it supports 
a variety of settings for the three key features of a Monte Carlo 
simulation: the move set, the acceptance criterion and the energy function. 

Some of the features available in Phaistos include:

\subsubsection*{Probabilistic models of protein structure}
\begin{itemize}
\item TorusDBN: a probabilistic model of the protein main chain \cite{boomsma2008gpm}.
\item BASILISK: a probabilistic model of amino acid side chains \cite{harder2010beyond}.
\item FB5HMM: a probabilistic model of a protein's $C\alpha$ trace \cite{hamelryck2006srp}.
\item COMPAS: a probabilistic model of the center of mass of amino acid side chains \cite{stovgaard2011}.
\end{itemize}
\subsubsection*{Move set}
\begin{itemize}
\item  Pivot, semilocal and local backbone moves in internal coordinate space.
\item  Pivot, semilocal and local backbone moves in internal coordinate space from probabilistic models that capture the local structure of proteins \cite{boomsma2008gpm}. 
\item  Rotamer moves, where tentative updates of one or more degrees of freedom are drawn from a uniform or Gaussian distribution.
\item  Rotamer moves from a probabilistic model that capture the side chain geometry of amino acids \cite{harder2010beyond}.
\end{itemize}
\subsubsection*{Simulation type}
\begin{itemize}
\item Metropolis-Hastings.
\item Greedy optimization.
\item Simulated annealing.
\item Generalized ensembles: multicanonical and $1/k$ \cite{frellsen12}.
\end{itemize}
\subsubsection*{Energy and likelihood functions}
\begin{itemize}
\item The OPLS-AA force field \cite{jorgensen1996development} with GB/SA implicit solvent model \cite{qiu1997gbsa}.
\item The Profasi force field \cite{irback2006profasi}.
\item Likelihood functions for SAXS \cite{stovgaard2010calculation, stovgaard2011} and NOE \cite{olsson2011generative} data.
\item Go-type potentials. 
\end{itemize}

\noindent In addition, a number of built-in analysis routines are
available (e.g fraction of native contacts, RMSD and helicity
calculation, cluster analysis, etc.).

A website containing an online version of this manual and the source
code of Phaistos can be found at
\emph{http://sourceforge.net/projects/phaistos/}. Bug reports,
questions and comments regarding Phaistos should be done through the
Sourceforge site.

%Phaistos is a software framework for the simulation of protein folding
%and the prediction of protein structure from sequence. It uses the
%Markov Chain Monte Carlo (MCMC) approach to molecular simulations, and
%includes a variety of efficient move-types to explore protein
%conformational space. Currently, Phaistos includes implementations of
%the two established forcefields: Profasi \cite{irback2006profasi} and
%OPLS \cite{jorgensen1996development} with the GBSA implicit solvent
%model \cite{qiu1997gbsa}.
%

\section{Citing Phaistos}
\label{sec:citing}

The main reference for Phaistos is:

\vspace{1em}\noindent W. Boomsma, J. Frellsen, T. Harder, S. Bottaro,
K. E. Johansson, P. Tian, K. Stovgaard, C. Andreetta, S. Olsson,
J. Valentin, L. D. Antonov, A. S. Christensen, M. Borg, J. H. Jensen,
K. Lindorff-Larsen, J. Ferkinghoff-Borg, T. Hamelryck,
J. Comput. Chem. 2013, DOI: 10.1002/jcc.23292.

\chapter{Installation}
\label{cha:installation}

Phaistos requires a C++ compiler, a fortran compiler, the boost
library\footnote{\url{www.boost.org}},
lapack\footnote{\url{www.netlib.org/lapack/}} and
CMake\footnote{\url{www.cmake.org}}. In addition, certain optional
modules may have additional requirements.

\section{Installation procedure}
\label{sec:inst-proc}

The recommended way to build Phaistos is out-of-source, meaning that the
generated binaries will not clutter the source directory. This is done
using the following sequence of commands:

\begin{verbatim}
$ mkdir build
$ cd build
$ cmake /path/to/phaistos/
$ make
\end{verbatim}

\noindent CMake can be configured from the command line using the -D syntax. For
instance, to specify a non-standard location of the boost directory, use:

\begin{verbatim}
$ cmake -DBOOST_ROOT:FILEPATH=directory ..
\end{verbatim}

\noindent The CMake output will contain a list of packages that were successfully
found and the corresponding CMake variables that can be modified if
necessary. A full list of the available CMake variables can be obtained
using:

\begin{verbatim}
$ cmake -LA ..
\end{verbatim}

\noindent By default, CMake will hide the details of the compilation
process. If you wish to see the executed commands, simply type:

\begin{verbatim}
$ make VERBOSE=1
\end{verbatim}

\section{Build type}
\label{sec:build-type}

\noindent Phaistos is by default built in release-mode, meaning that
the produced code is optimized for fast execution. During debugging,
it is often an advantage to compile to non-optimized code. This can be
done by modifying the BUILD\_TYPE variable:

\begin{verbatim}
$ cmake -DCMAKE_BUILD_TYPE:STRING=Debug ..
\end{verbatim}

\noindent If you are not satisfied with the default compiler settings,
you can always override them yourself:

\begin{verbatim}
$ cmake -DCMAKE_CXX_FLAGS="-ggdb -Wall" ..
\end{verbatim}

\section{Modules}
\label{sec:modules}

Phaistos has a module system, which makes it easy to write your own
code (module) and compile it in with the main library. Parts of the
Phaistos functionality is provided as modules, to allow the user to
exclude it during compilation, which can speed up compilation times
significantly. CMake will output a list of all detected modules:

\begin{verbatim}
-- Modules - phase: 1
--   Found modules:
--     git
--     muninn
--     pleiades
...
\end{verbatim}

\noindent
You can disable modules using the \texttt{PHAISTOS\_MODULE\_DISABLE}
and \\ \texttt{PHAISTOS\_MODULE\_ENABLE} options. Both take regular
expressions specifying which modules to disable or enable,
respectively. For instance, if we wish to include only the \texttt{git}
module, this can be done by enabling it directly

\begin{verbatim}
$ cmake -DPHAISTOS_MODULE_ENABLE:STRING="git" ..
\end{verbatim}

\noindent
or by disabling everything else

\begin{verbatim}
$ cmake -DPHAISTOS_MODULE_DISABLE:STRING="muninn|pleiades|..." ..
\end{verbatim}

\noindent
In the list of detected modules, CMake will indicate which modules are
currently disables:

\begin{verbatim}
-- Modules - phase: 1
--   Found modules:
--     git
--     muninn (disabled)
--     pleiades (disabled)
\end{verbatim}

CMake variables are cached, so if you specify an option once it will
be remembered in future executions of CMake. Therefore, remember to
explicitly remove the \texttt{PHAISTOS\_MODULE\_DISABLE} or
\texttt{PHAISTOS\_MODULE\_ENABLE} option value if you no longer need
it. This can be done by setting it to the empty string:

\begin{verbatim}
$ cmake -DPHAISTOS_MODULE_DISABLE:STRING="" ..
\end{verbatim}

\section{Constructing a source tar-ball}
\label{sec:constr-source-tar}

You can create a tar-ball of all Phaistos source files by using the
following command:

\begin{verbatim}
$ make package_source
\end{verbatim}

\noindent
This will create a file called \texttt{phaistos-XXX.tar.gz} where
\texttt{XXX}, is the current version number.

\section{Locating data files (\texttt{PHAISTOS\_ROOT})}
\label{sec:locating-data-files}

Phaistos uses a number of data files, which are typically referred to
using a path relative to the execution directory. This means that
\texttt{phaistos.cpp} will be able to find these files when started
from the \texttt{bin} directory in your build directory, but it might
fail if you try to run it from another directory. This problem can be
resolved by setting the \texttt{PHAISTOS\_ROOT} environment variable
to point to the root of your build directory. In the bash shell, this
can be done using:

\begin{verbatim}
$ export PHAISTOS_ROOT='/path/to/phaistos/build/'
\end{verbatim}

\noindent
while tcsh requires a slightly different syntax:

\begin{verbatim}
> setenv PHAISTOS_ROOT '/path/to/phaistos/build/'
\end{verbatim}



\chapter{Running simulations using \texttt{phaistos.cpp}}
\label{cha:running-simulations}

The main executable in the Phaistos package is called
\texttt{phaistos.cpp}. This program is a front-end to basically all
functionality in the library. Settings for \texttt{phaistos.cpp} are
specified using either command line options or a configuration
file. Most settings in Phaistos have default values, so typically, a
simulation can be set up by specifying only a few parameters. As the
simulation starts, it will output all used settings to the screen,
allowing the user to verify that the current configuration is
meaningful. This output is formatted so that it is compatible with the
configuration file format, meaning that changes to the current setting
can be made simply by copy-and-pasting from the initial output to a
configuration file.

\section{Command line options}
\label{sec:command-line-options}

\texttt{phaistos.cpp} uses long unix-style command line options of the form:

\begin{flushleft}
\texttt{--option-name \emph{optional-value}}. 
\end{flushleft}

\noindent All available options for \texttt{phaistos.cpp}, along with
their default values, can be obtained by using with the help options:

\begin{verbatim}
$ ./phaistos --help
\end{verbatim}

\noindent For convenience, \texttt{phaistos.cpp} has shorthand
options, which will automatically expand to their longer
equivalents. The most important examples are the \texttt{move} and
\texttt{energy} shorthands. These allow you to quickly specify sets of
moves and energy terms. For instance

\begin{verbatim}
$ ./phaistos --energy clash-fast opls-angle-bend opls-vdw
\end{verbatim}

\noindent
corresponds to

\begin{verbatim}
$ ./phaistos --energy-clash-fast --energy-opls-angle-bend \ 
  --energy-opls-vdw
\end{verbatim}

\noindent The individual energy terms might have options of their
own. These can be specified in square brackets in the shorthand
notation:

\begin{verbatim}
$ ./phaistos --energy clash-fast[debug:1,weight:10]
\end{verbatim}

\noindent
This will expand to

\begin{verbatim}
$ ./phaistos --energy-clash-fast --energy-clash-fast-debug 1 \ 
  --energy-clash-fast-weight 10
\end{verbatim}

\subsection{Procedures}
\label{sec:procedures}

The \texttt{phaistos.cpp} executable contains several main functions,
which address different types of simulation problems. These main
functions are called \emph{procedures}, and can be selected from the
command line using the \texttt{--procedure} option. The default
procedure is \texttt{fold}, which runs a general protein
simulation. There is a procedure called \texttt{compactify}, which
rapidly produces compact structures, and debug procedure called
\texttt{comparison}. The latter allows the user to easily compare results
of two energy functions, for instance to verify that a cached version
of an energy produces the same result as the non-cached version.

Procedures have sub-options which are relevant only to that particular
procedure. For instance, the \texttt{fold} procedure has options to
control whether debug information should be written out to screen
during a simulation. This can be specified using the
nested style introduced in the previous section:

\begin{verbatim}
$ ./phaistos --procedure fold[debug:true]
\end{verbatim}


\subsection{Modes}
\label{sec:modes}

Modes are shortcuts that associate a name to a combination of specific
simulation settings that are often used. Currently, there are only a
few predefined modes, but more will be added in the future. A mode is
selected using the \texttt{--mode} command-line option. One example is
the \texttt{opls-mc-dynamics} mode, which will be illustrated in an
example below (see section \ref{sec:study-native-ensembl}).


\section{Configuration files}
\label{sec:configuration-files}

As the \texttt{phaistos.cpp} program is started, it will output a complete list of
settings, including both the ones you specified from the command line,
and default values for all other settings. This output, starting with

\begin{verbatim}
##################### PHAISTOS OPTIONS #####################
\end{verbatim} 

\noindent
and ending in

\begin{verbatim}
################### PHAISTOS OPTIONS END ###################
\end{verbatim}

\noindent
is in a format that \texttt{phaistos.cpp} can read in as a configuration file. If you
copy-and-paste these options into a file \texttt{filename}, and run Phaistos using

\begin{verbatim}
$ ./phaistos --config-file filename
\end{verbatim}

\noindent
it will use the same settings as those provided in the command line
originally. This gives a convenient workflow for setting up a
simulation: start by providing a few settings for the command line,
copy-and-paste the resulting output to a configuration file, and then
carefully read through (and modify) this configuration file to ensure
that you have the desired settings.


\section{Examples}
\label{sec:examples-simulations}

We include a few examples of possible simulation configurations. This
is mainly to give a broad overview of how simulations are set up. The
move, energy, and simulation settings will be covered in detail in the
next chapters.

\subsection{Studying a native ensemble using OPLS}
\label{sec:study-native-ensembl}

The \texttt{opls-mc-dynamics} mode provides a simple example of a
simulation using the OPLS force field:

\begin{verbatim}
$ ./phaistos --mode opls-mc-dynamics --pdb-file 2gb1.pdb \
  --init-from-pdb
\end{verbatim}

\noindent
We explicitly state which PDB file we are using, and that you wish to
initialize the simulation from this PDB file. The rest of the settings
are given by the \texttt{opls-mc-dynamics} mode. This mode expands to
settings for the complete OPLS forcefield, and a variety of moves that
were found to be optimal in the context of exploring native ensembles
\cite{bottaro2012crisp}. In addition, the simulation type is set to
Metropolis-Hastings. Note that both the \texttt{crisp} and \texttt{opls} modules must be
enabled for this mode to work (which they are by default).

A mode simply provides default values for certain settings. Everything
can be overridden by specifying additional options from the command
line. For instance, you can adjust the number of iterations and dump
pdb files at certain intervals by adding the following options to the
command line:

\begin{verbatim}
--iterations 100000000 --observable pdb[register-interval:20000]
\end{verbatim}


\subsection{Generalized ensemble simulation using Profasi}
 
Profasi is a highly efficient forcefield, that has been used in the
literature both for folding and aggregation simulations
\cite{irback2009effective}. This, in combination with the Muninn
generalized ensemble framework \cite{frellsen12}, is a powerful tool for conducting
rapid equilibrium folding simulations.

A simple version of such a simulation could look like this:

\begin{verbatim}
./phaistos --aa-file 1L2Y.aa --move pivot-uniform sidechain-uniform \
 --energy profasi-cached --threads 2 \
 --monte-carlo muninn[min-beta:0.05,max-beta:1.25]
\end{verbatim}

\noindent
In this case, we start from an amino acid sequence file, containing
just a single line of amino acids:

\begin{verbatim}
NLYIQWLKDGGPSSGRPPPS
\end{verbatim}

\noindent
In addition, we specify that we wish to use uniform moves both for the
protein backbone and sidechain, and that the simulation is done using
the cached version of the profasi forcefield. We use the Muninn
generalized ensemble framework for this simulation, and set up the
temperature range (\texttt{min-beta} and \texttt{max-beta}). Finally,
the \texttt{threads} option specifies that we wish to run with two
threads. This has the effect that we will be conducting two
independent simulations, which uses the same set of weights and both
will contribute to the same histograms and entropy estimates in
Muninn.


\chapter{Moves}
\label{sec:moves}

Phaistos includes a range of different approaches to update the
molecule chain in each iteration of the simulation. These are referred
to as \emph{moves}. Unless explicitly mentioned, all moves in Phaistos
obey detailed balance. This means that if you simulate with for
instance standard Metropolis-Hastings and no energy function, you can
expect flat distributions over the dihedral angle degrees of freedom
of your system. More generally, if you simulate with
Metropolis-Hastings and a given \emph{physical} energy function, you
have the Boltzmann distribution as your target distribution, see
section \ref{sec:energy_convention} for details.

\section{General move settings}

All moves share the following basic settings. The \texttt{weight}
setting determines how often the move is used in the simulation, while
the \texttt{regions} option allows you to specify where in the chain
you want to apply the move. The \texttt{debug} flag can be used (if
implemented by the author of the move) to output additional
information about the move.

\begin{optiontable}
\option{weight}{real}{1.0}{Relative weight of the move (probability of choosing this move).}
\option{debug}{int}{0}{Debug level.}
\option{regions}{vector$<$pair$<$int$>>$}{[[0,chain\_size]]}{Regions affected by move.}
\end{optiontable}

\section{Uniform pivot move (\texttt{pivot-uniform})}
\label{sec:pivot-move-uniform}

Perhaps the most classic Monte Carlo move in molecular simulations,
the pivot-uniform simply produces a random, uniformly distributed
rotation of the dihedral angles ($\phi$,$\psi$ angles) in a single
residue. The \texttt{max-delta} option determines the maximum size of the
rotation. You can use the \texttt{single-dof-only} option to limit yourself to
a single degree of freedom. The \texttt{skip-proline-phi} option makes it
possible to omit the $\phi$ dihedral in Prolines, which involves a
change in bond length and should be avoided if no bond length term is
included in the energy.

\optiontitle{Settings}
\begin{optiontable}
  \option{move-length-min}{int}{1}{Minimum length of move.}
  \option{move-length-max}{int}{1}{Maximum length of move.}
  \option{single-dof-only}{bool}{true}{Whether to update only one $\phi$ or $\psi$ angle at a time (randomly chosen).}
  \option{skip-proline-phi}{bool}{true}{Whether to sample angles in proline. Note, if \texttt{skip-proline-phi} is set to false, a change in proline bond length is introduced, which must be taken in account by an appropriate bond-stretch energy term.}
 \option{max-delta}{float}{$\pi$(rad)}{Maximum change in angle value.}
 \end{optiontable}

\section{Backbone DBN move (\texttt{backbone-dbn})}
\label{sec:pivot-move-dbn}

The Backbone DBN move is essentially a pivot move using a non-trivial
probability distribution. More precisely, angular degrees of freedom
are sampled from a probabilistic model of the protein
backbone. Currently, two models are available: TorusDBN and Fb5Dbn
\cite{boomsma2008gpm,hamelryck2006srp}, using the full-atom and the
C$_{\alpha}$-only representation respectively. The relevant model is
selected using the \texttt{backbone-dbn} option (see chapter
\ref{cha:probabilistic-models}). When using TorusDBN, the move
basically resamples ($\phi$,$\psi$) dihedral angles from a
Ramachandran like distribution.

\notification{IMPORTANT: This move introduces a bias (implicit energy) to the simulation that is possible to compensate by setting \texttt{implicit-energy}=false.}
%\notification{IMPORTANT: This move introduces an implicit energy to the simulation (unless \texttt{implicit-energy}=false).}
\optiontitle{Settings}
\begin{optiontable}
  \option{move-length-min}{int}{1}{Minimum length of move.}
  \option{move-length-max}{int}{7}{Maximum length of move.}
  \option{resample-mode}{Enum}{RESAMPLE\_ALL}{Which values to resample in the model.}
  \option{implicit-energy}{bool}{true}{Whether the bias introduced by the dbn move should be included as an implicit energy.}
  \option{dbn-consistency-window-size}{int}{10}{Whenever a dbn move is combined with a different move, the dbn will need to be made consistent with the current state of the protein after every non-dbn move. Ideally, this requires a full-length resampling (corresponding to option value -1). However, since the dbn has a finite memory, a window of 10 is a reasonable approximation.}
  \option{dbn-bias-window-size}{int}{10}{To evaluate the bias introduced by a dbn move, ideally, a full-length forward calculation must be done. However, since the dbn has a finite memory, a window of 10 is a reasonable approximation.}
\end{optiontable}


\section{Local pivot move (\texttt{pivot-local})}
\label{sec:pivot-move-gaussian}

The \texttt{pivot-local} move modifies $(\phi, \psi)$ angles and
optionally non-standard degrees of freedom (e.g.\ bond angles or
$\omega$ angles) in the backbone. Updates are proposed from a Gaussian
distribution centered on the current value and with user-controlled
variance.

% WHY IS THE INTERFACE OF THIS MOVE DIFFERENT FROM THE UNIFORM VARIANT?

\optiontitle{Settings}
\begin{optiontable}
  \option{move-length-min}{int}{1}{Minimum length of move.}
  \option{move-length-max}{int}{1}{Maximum length of move.}
  \option{sample-phi-psi-dofs}{bool}{true}{Whether to modify $(\phi, \psi)$ angles.}
  \option{sample-bond-angle-dofs}{bool}{true}{Whether to modify bond angles.}
  \option{sample-omega-dofs}{bool}{false}{Whether to modify omega angles.}
  \option{std-dev-phi-psi}{real (degrees)}{4}{Parameter controlling
    the variance of the Gaussian distribution for $(\phi, \psi)$ angles (degrees).}
  \option{std-dev-bond-angle}{real (degrees)}{0.8}{Parameter controlling
    the variance of the Gaussian distribution for bond angles (degrees).}  
  \option{std-dev-omega}{real (degrees)}{0.8}{Parameter controlling
    the variance of the Gaussian distribution for $\omega$ angles (degrees).}
\end{optiontable}

There is also a version of this move that includes an implicit DBN
energy for dihedral angles and an implicit bond angle energy based on
Engh-Huber statistics. This variant is called
\texttt{pivot-local-dbn-eh}.


\section{Uniform sidechain move (\texttt{sidechain-uniform})}
\label{sec:sidechain-uniform-move}

The \texttt{sidechain-uniform} move produces random, uniformly distributed
rotations of the sidechain torsion angles ($\chi$ angles) in single
residues.

\optiontitle{Settings}
\begin{optiontable}
  \option{move-length-min}{int}{1}{Minimum length of move.}
  \option{move-length-max}{int}{1}{Maximum length of move.}
  \option{single-dof-only}{bool}{true}{Whether to update only one $\chi$ angle at a time.}
  \option{skip-proline}{bool}{true}{Whether to sample sidechain angles in proline. Note, if \texttt{skip-proline} is set to false a change in proline bond length is introduced, which must be taken in account by an appropriate bond-stretch energy term.}
 \end{optiontable}



\section{Rotamer sidechain move (\texttt{sidechain-rotamer})}
\label{sec:sidechain-rotamer-move}

The rotamer move (\texttt{sidechain-rotamer}) modifies the $\chi$-angles in
the sidechain. New angles are proposed from the Gaussian distributions
in the Dunbrack backbone-independent rotamer library
\cite{dunbrack1997bayesian}.

\notification{IMPORTANT: This move introduces a bias (implicit energy) to the simulation that is possible to compensate by setting \texttt{implicit-energy}=false.}
%\notification{IMPORTANT: This move introduces an implicit energy to the simulation (unless \texttt{implicit-energy}=false).}

\optiontitle{Settings}
\begin{optiontable}
  \option{move-length-min}{int}{1}{Minimum length of move.}

  \option{move-length-max}{int}{1}{Maximum length of move.}

  \option{implicit-energy}{bool}{true}{If
  implicit-energy is set to false, the bias introduced when
  proposing from the Dunbrack library is compensated for, ensuring
  that the uniform distribution is reached when
  simulating with no force-field.}

  \option{--sigma-scale-factor}{real}{1.0}{The width of the Dunbrack
  Gaussian distributions from which new angles are proposed can be
  scaled by this factor.  Flattening the distribution can be useful to
  ensure ergodicity, as the Dunbrack library does not have support
  over the entire interval $[0,2\pi[$.}

  \option{--rotamer-state-resample-frequency}{real}{1.0}{Frequency
  with which move will resample new rotamer states ($g^+, t, g^-$).}

  \option{--rotamer-skip-proline}{bool}{true}{Whether to sample
  sidechain angles in proline. Note, if rotamer-skip-proline is set to false a
  change in proline bond length is introduced, which must be taken in
  account by an appropriate bond-stretch energy term.}

  \option{--sample-hydrogen-chis}{bool}{true}{Whether to sample
  (uniformly) non-standard torsion angles in sidechain (e.g methyl
  groups, CB-S-H in Cys, etc...).}
\end{optiontable}


\section{Local sidechain move (\texttt{sidechain-local})}
\label{sec:sidechain-local-move}

The \texttt{sidechain-local} move modifies $\chi$ angles and optionally
non-standard degrees of freedom (e.g. bond angles) in the sidechain.
Four different strategies can be employed:

\begin{description}
\item[simple] Tentative updates are proposed for all the degrees of
  freedom in the sidechain from a Gaussian distribution centered on
  the current angle(s) and with user-controlled variance
\item[scale-sigma] Tentative updates are proposed for all the degrees
  of freedom in the sidechain from a Gaussian distribution centered on
  the current angle(s).  The width of the Gaussians for the different
  angles in the sidechain is scaled linearly so that degrees of
  freedom further away from CA have larger variance.
\item[select-dofs] Tentative updates are proposed for one degree of
  freedom in the sidechain from a Gaussian distribution centered on
  the current angle and with user-controlled variance.  The degree of
  freedom is chosen at random according to a weight given by its
  distance to the CA atoms, such that angles closer to the backbone
  are sampled less frequently.
\item[constrain-endpoint] Tentative updates are proposed for all the
  degrees of freedom in the sidechain from a multivariate Gaussian
  distribution with a bias towards small displacement of atoms
  involved as acceptors or donors in hydrogen bonds.  This type of
  move is is useful to enable small adjustment of the sidechain
  without breaking the dense network of non-covalent interactions.
\end{description}

\optiontitle{Settings}
\begin{optiontable}
  \option{move-length-min}{int}{1}{Minimum length of move.}  
  \option{move-length-max}{int}{1}{Maximum length of move.}
  \option{sample-major-dofs}{bool}{true}{Whether to sample $\chi$ angles.}
  \option{sigma-major-dofs}{real(radians)}{0.03490}{Width (in radians) of the Gaussian distribution for $\chi$ angles.}
  \option{sample-minor-dofs}{bool}{true}{Whether to sample minor degrees of freedom (all bond-angles and non-$\chi$ dihedral angles).}
  \option{sigma-minor-dofs}{real}{0.00349}{Width (in radians) of the Gaussian distribution for minor dofs angles.}
  \option{mode}{string}{simple}{Four different move modes are available: simple, scale-sigma. select-dofs, constrain endpoint (described above).}
  \option{rotamer-lagrange-multiplier}{real}{200}{When using constrain-endpoint mode, this parameter controls the weight of the constraint.}
  \option{skip-proline}{bool}{true}{Whether to sample sidechain angles in proline. Note, if \texttt{skip-proline} is set to false a change in proline bond length is introduced, which must be taken in account by an appropriate bond-stretch energy term.}
\end{optiontable}


\section{Crankshaft move (\texttt{crankshaft})}

The CRANKSHAFT move consists of a rigid rotation of a chain segment
about the axis defined by the $\mathrm{C}_{\alpha}$ atoms delimiting
the segment \cite{betancourt2005efficient}. This move therefore alters
only the backbone dihedral angles $\phi$ and $\psi$ and the
$\mathrm{C}_{\alpha}$ bond angles of the segment ends. The move is
implemented as described in reference \cite{smith2008backrub}.  The
original procedure includes an optimized placement of
$\mathrm{C}_{\beta}$ and $\mathrm{H}_{\alpha}$ atoms, which does not
fulfill detailed balance, which is therefore omitted in this
implementation.


\optiontitle{Settings}
\begin{optiontable}
  \option{move-length-min}{int}{2}{Minimum length of move.}  
  \option{move-length-max}{int}{12}{Maximum length of move. The number of residues involved in a single update is randomly chosen in the interval [\texttt{move-length-min},\texttt{move-length-max}].}
  \option{constrain-bond-angle}{bool}{true}{Whether to apply a constraint on the bond angles variation.}
  \option{bond-angle-tolerance}{real (degrees)}{15}{If constrain-bond-angle is true, this parameter sets the maximum variation in bond angles (in degrees) from the original structure.}
  \option{only-internal-moves}{bool}{false}{Whether to perform only internal moves or both internal and end-moves. End-moves consist is small variations of bond- and torsion angles.}
\end{optiontable}



\section{Sidechain DBN moves}
\label{sec:sidechain-move-dbn}

The sidechain DBN moves (\texttt{sidechain-basilisk} and \texttt{sidechain-compas})
allow you to sample amino acid side chain conformations from one of
the included probabilistic models of side chain geometry, namely
COMPAS or BASILISK \cite{harder2010beyond}.  These moves will propose
a new set of angles for a randomly selected stretch of residues. We
recommend only using a stretch of length one, or in other words, to
only resample a single side chain at a time. Changing multiple
sidechains in one go will increase the likelihood of introducing
clashes, leading to a significantly higher rejection rate.

Both moves use an external model file. Please remember to set the
\texttt{PHAISTOS\_ROOT} environmental variable to avoid problems
accessing these files (see section \ref{sec:locating-data-files}).

\subsection{General settings}

The behavior of the two moves is identical in terms of
parameters. However, the moves applies to different side chain
representation: BASILISK requires full atomic details, whereas COMPAS
only requires a pseudo atom at the side chain center of mass.  Both
sidechain DBN moves have the following options.

\begin{optiontable}
\option{move-length-min}{int}{1.}{Minimum move length.}
\option{move-length-max}{int}{1.}{Maximum move length.}
\option{mocapy-dbn-dir}{string}{../data/mocapy\_dbns}{Path to mocapy model file directory.}
\option{implicit-energy}{bool}{\emph{true}}{Whether the dbn bias (implicit energy) should be divided out (=false) or not (=true).}

\option{ignore-bb}{bool}{\emph{false}}{Whether to use the backbone independent version instead.}
\option{check-move}{bool}{\emph{true}}{Check that the move produced consisted atom positions.}
\option{sample-hydrogen-chis}{bool}{\emph{false}}{Sample sidechain residual dofs.}
\option{sample-hydrogen-chis-normal}{real}{1.}{Resampling constrained to a Gaussian around the current state.}
\option{sample-hydrogen-chis-sigma}{real}{0.035}{Standard deviation of Gaussian for resampling residual dofs.}  
\option{model-filename}{string}{basilisk.dbn/compas.dbn}{Model filename.}
\end{optiontable}

\notification{IMPORTANT: Both variants of the move introduce a bias (implicit energy) to the 
simulation that is possible to compensate by setting \texttt{implicit-energy}=false.}

\subsection{BASILISK move (\texttt{sidechain-basilisk})}
\label{sec:basilisk-move}

BASILISK assumes a full atom representation of the protein
sidechain. There are several variants of this move available:

\begin{flushleft}
\texttt{sidechain-basilisk}, \texttt{sidechain-basilisk-local}, and \texttt{sidechain-basilisk-multi}.
\end{flushleft}

\noindent The \texttt{local} version samples the hidden node state
based on the current angle values and subsequently samples angles
corresponding to this state, which generally has the effect of small
angle variations. The \texttt{multi} variant will do a sidechain move
in two disjunct regions of the protein chain.  See appendix
\ref{cha:sidechain-dbns-module} for details.


\subsection{COMPAS move (\texttt{sidechain-compas})}
\label{sec:compas-move}

COMPAS uses a pseudo-sidechain representation, with one atom per
sidechain. See appendix \ref{cha:sidechain-dbns-module} for details.



\section{None move (\texttt{none})}
\label{sec:none-move}

The \texttt{none} move does no structural resampling. It is included
because some energy terms are constructed to resample their own
parameters, in which case it can be convenient to have a Monte Carlo
iteration where the structure is not modified.

\section{Selecting moves in \texttt{phaistos.cpp}}

Moves can be selected from the \texttt{phaistos.cpp} command line
either individually

\begin{verbatim}
$ ./phaistos --move-pivot-uniform --move-sidechain-uniform
\end{verbatim}

\noindent
or using a shorthand

\begin{verbatim}
$ ./phaistos --move pivot-uniform sidechain-uniform
\end{verbatim}

\noindent
When using the shorthand option, indivudial move options can be
specified using square brackets as described in section
\ref{sec:command-line-options}.



\chapter{Energies}
\label{sec:energies}

\section{Phaistos energy convention}
\label{sec:energy_convention}

To ensure consistency between probabilistic terms and physical energy
terms, energies in Phaistos, $E_{\phi}$, are reported as the negative
logarithm of the probability of observing a structure. This means that

\begin{equation}
E_{\phi}(x) = -\log P(x) \ ,
\end{equation}

\noindent where $P(x)$ is the probability of observing the structure
$x$ according to the energy term, and $E_{\phi}(x)$ is the
corresponding Phaistos energy. As a consequence Phaistos energies are
unitless.

For classical physics based energy terms, this means that the Phaistos
energy becomes $E_{\phi}(x) \propto \beta {E(x)}$, where $E$ the physical
energy function and $\beta = (kT)^{-1}$ is the thermodynamic
beta. Here, $k$ is the Boltzmann constant and $T$ is the
temperature. The $\beta$ weight is specified using the \texttt{weight}
setting of the individual energy terms. Physical energies will often
be carrying a default weight of $1.67857$ at $300\,\mathrm{K}$ if the
physical energy is measured in $\mathrm{kcal}/\mathrm{mol}$. This
weight corresponds to the value of $(kT)^{-1}$ for $T=300\,\mathrm{K}$
and $k = 1.9872159 \cdot 10^{-3}
\frac{\mathrm{kcal}}{\mathrm{mol}\,\mathrm{K}}$.


\section{General energy settings}
\label{sec:gener-energy-sett}

The following settings are shared by all energy terms. The
\texttt{weight} is multiplied onto the energy by the energy term
itself and consequently before the summation of energy terms when
calculating the total energy of a structure. Note, the energy values
outputted during a simulation is alway multiplied by the weight
associated with the energy term. For physical energy terms, the weight
corresponds to usual $1/kT$ factor (see previous section). Just as for
the moves, all energies also have a \texttt{debug} flag, which
indicates that the energy evaluation should execute in debug mode.

\begin{optiontable}
\option{weight}{real}{1.0}{Absolute weight used when evaluating the weighted energy used for summing energy terms.}
\option{debug}{int}{0}{Debug level.}
\end{optiontable}

\section{Clash detection (\texttt{clash-fast})}
\label{sec:clash-fast}

Phaistos employs a ChainTree data structure to quickly determine
interacting atoms within a given distance
\cite{lotan2004algorithm}. The \texttt{clash-fast} energy function
uses this data structure to do a coarse-grained self-collision check, returning
infinity if atoms are found to be within the given cutoffs. This makes
it useful as a first energy term in an energy function, acting as a
quick filter. This works because remaining energy terms are skipped if
the energy becomes infinite during the evaluation.

\optiontitle{Settings}
\begin{optiontable}
  \option{only-modified-pairs}{bool}{true}{Specifies whether the energy should consider only pairs modified by the last move. If a clashfree structure is maintained at all times, this is sufficient to detect all clashes.}
  \option{boolean-mode}{bool}{true}{Specifies whether the energy should work in clash/non-clash mode and return infinity/0 (true) or count all the clashes and return the number of clashes (false).}
  \option{minimum-residue-distance}{int}{2}{Minimum distance along chain (measured in number of residues) before pair is taken into account by the energy (in angstrom).}
  \option{clash-distance-h}{real}{1.5}{Distance within which pairs of hydrogens are considered to be clashing (in angstrom).}
  \option{clash-distance-no}{real}{2.3}{Distance within which pairs of Nitrogen-Oxygen pairs are considered to be clashing (in angstrom).}
  \option{clash-distance-ps}{real}{1.9}{Distance within which pairs of pseudo-sidechain atom pairs are considered to be clashing  (in angstrom).}
  \option{clash-distance-sg}{real}{1.8}{Distance within which pairs of SG atom pairs are considered to be clashing (in angstrom).}
  \option{clash-distance-any-pair}{real}{2.3}{Default distance within which pairs of atom pairs of all other types are considered to be clashing (in angstrom).}
\end{optiontable}

\section{Backbone DBN energy (\texttt{backbone-dbn})}
\label{sec:backbone-dbn-energy}

This is a backbone energy that uses one of the probabilistic models
described in chapter \ref{cha:probabilistic-models}. The model itself
is selected using the \texttt{backbone-dbn} option (see chapter
\ref{cha:probabilistic-models}).

In principle, when another move modifies the backbone angles of a
segment of the protein backbone, this affects the distribution over
all hidden nodes in the sequence according to the probabilistic
model. In practice, we know that the models have a finite memory, and
it is sufficient to recalculate only a window around the changed
segment. The size of this window is controlled using the
\texttt{window-size} parameter.

\optiontitle{Settings}
\begin{optiontable}
  \option{enable-dbn-update}{bool}{true}{Whether to update the angles in the DBN from the chain when necessary.}
  \option{always-full-update}{bool}{false}{Whether to always force update of all angles in the DBN from the chain.}
  \option{window-size}{int}{-1}{Size of window used when calculating the energy. A good value for the window size is $>7$, and a negative window size means that the full bias is used.}
\end{optiontable}



\section{MuMu: Non-local multibody energy}
The multibody multinomial (MuMu) model is a coarse grained probabilistic model of non-local interactions in proteins \cite{johansson2013mumu}. 

The main usage of this module is for residue probability annotation of PDB files. This is easily done with the following line:
\begin{flushleft}
\texttt{./evaluate\_observable --pdb-file xxxx.pdb --observable-mumu output-target=pdb-b-factor}
\end{flushleft}
Note that the current implementation does not check for missing atoms or modified residues. The module will use the atoms provided by the Phaistos PDB parser. Missing or unknown atoms will result in erroneous evaluations.

The usage as a potential of mean force in a folding simulation will require a reference model which is not provided. However, with a weight of 2-5 the energy will stabilize most structures.

\optiontitle{Settings}
\begin{optiontable}
  \option{use-ratio}{bool}{0}{Use the ratio model. This is not available.}
  \option{model-filename}{string}{default}{Name of the Mocapy *.dbn model file. If set to 'default' the file name will be read from mumu\_data\_miner.h}
  \option{reference-model-filename}{string}{default}{Name of the Mocapy *.dbn reference model file. If set to 'default' the file name will be read from mumu\_data\_miner.h. Not provided.}
\end{optiontable}



\section{Contact-map energy (\texttt{contact-map})}
\label{sec:contact-map}

This is an artificial energy term that can be used to drive proteins
to the correct fold by providing native contact information. There are
two supported functional forms: $w \exp\frac{-(d-d_0)^2}{\sigma^2}$
(default) and $w \frac{(d-d_0)^2}{\sigma^2}$ (selected by setting
\texttt{dist-squared}=false), where $\sigma$ is called the contact
width and $w$ is the contact weight. The default value for the contact
width can be set by the option \texttt{default-width}, whereas the
default value for the contact weight is 1. The contact map can be
defined with a contact map file. The input file takes one contact per
line, in the following format:

{
\footnotesize
\begin{verbatim}
residueno1 atomname1 residueno2 atomname2 idealdistance [width [weight]]
\end{verbatim}
}

\noindent In this format, the width and the weight is optional. The
input format could for example look like:

\begin{verbatim}
5 CA 1 CA 13.029
6 CA 2 CA 13.168 2.0 1.5
\end{verbatim}

\optiontitle{Settings}
\begin{optiontable}
  \option{contact-map-file}{string}{}{Contact map input file.}
  \option{contact-map-string}{string}{}{Contact map input string.}
  \option{pdb-file}{string}{}{Path to PDB file for constructing contact map.}
  \option{cutoff}{real}{inf}{contact cutoff distance when constructing map from PDB.}
  \option{minimum-residue-distance}{int}{7}{Minimum residue-separation in a pair.}
  \option{iteration-type}{IterateEnum}{ALL}{Which atoms to include when constructing contactmap from chain.}
  \option{dist-squared}{bool}{false}{Whether to use the squared distance version of the energy.}
  \option{potential-width}{real}{1.0}{Width of energy potential.}
\end{optiontable}

\section{Selecting energies in \texttt{phaistos.cpp}}

In the \texttt{phaistos.cpp} command line, energy terms can either be
specified one by one

\begin{verbatim}
$ ./phaistos --energy-clash-fast --energy-opls-angle-bend
\end{verbatim}

\noindent
or using a shorthand

\begin{verbatim}
$ ./phaistos --energy clash-fast opls-angle-bend
\end{verbatim}

\noindent
When using the shorthand option, indivudial energy options can be
specified using square brackets as described in section
\ref{sec:command-line-options}.



\chapter{Observables}
\label{sec:observables}

An observable is an energy term with some additional functionality,
allowing the user to specify how the value should be outputted and how
frequently. All energy terms can be used as observables, and
conversely, all observables are based on an underlying energy term.


\section{General observable settings}

In addition to the general settings of an energy term (see
section~\ref{sec:gener-energy-sett}), all observables have the
following settings.

\begin{optiontable}
\option{register-interval}{int}{5000}{Interval (number of iterations) with which observable should be registered.}
\option{register-burnin}{int}{0}{Interval (number of iterations) before observable should start registering.}
\option{output-target}{string}{observables\_\%p\_\%t.dat}{Where and how the observable should be dumped (see text).}
\option{output-interval}{int}{1}{Interval (number of iterations) with which observable should be outputted. This value is ignored if it is lower than register-interval.}
\end{optiontable}

\noindent
The \texttt{output-target} setting can take various predefined 
values.

\begin{description}
\item{\texttt{stdout}$|$\texttt{cout}} Output to \texttt{stdout}.
\item{\texttt{stderr}$|$\texttt{cerr}} Output to \texttt{stderr}.
\item{\texttt{pdb-header}} Output to the header in dumped PDB files.
\item{\texttt{pdb-b-factor}} Output to the b-factors in dumped PDB files.
\end{description}

\noindent
Any other values will be regarded as a filename specification. The
strings \texttt{\%p} and \texttt{\%t} will be replaced by the current
process id and thread number, respectively. If a \texttt{\%i} substring
is found in the filename, it will be replaced by the current iteration
of the simulation (meaning that there will be a new log-file for each
iteration in which an observation is made).

It is possible to specify several output targets for one
observable. This is done by separating them with a \texttt{+}. For
instance, this command line will output both to a logfile and to
stdout:

\begin{verbatim}
$ ./phaistos --observable \
  helix-content[output-target:stdout+logfile.dat]
\end{verbatim}

\subsection{Output modes}
\label{sec:output-modes}

There are a number of different output modes available:

\begin{flushleft}
\texttt{verbose}, \texttt{verbose-with-id}, \texttt{verbose-with-time}, \texttt{compact}, \texttt{minimal}.
\end{flushleft}

\noindent These are normally chosen automatically based on the
output-target. However, they can be overridden by providing an
\texttt{\#\emph{output-mode}} in an output-target. For instance

\begin{verbatim}
$ ./phaistos --observable \
  helix-content[output-target:logfile.dat#verbose]
\end{verbatim}

\noindent
will force the observable to be written to the \texttt{logfile.dat}
file in \texttt{verbose} mode.

\section{Radius of gyration (\texttt{rg})}
\label{sec:radius-gyration}

The radius of gyration energy term is mostly used as an observable,
and is therefore described in this section, although it could equally
well have been included in chapter \ref{sec:energies}. It calculates
the root mean square distance of the atoms to the center of mass of
the protein.

\section{Root mean square deviation (\texttt{rmsd})}
\label{sec:root-mean-square}

This is another energy term that is mostly used as an observable, but
can serve as an artificial energy function as well. It returns the
Root Mean Square Deviation between the current structure and a
reference structure.

\optiontitle{Settings}
\begin{optiontable}
\option{reference-pdb-file}{string}{}{Path to a PDB file to which comparison should be done.}
\option{ca-only}{bool}{true}{Whether only to include CA atoms in the calculation.}
\end{optiontable}

\section{Helix content (\texttt{helix-content})}
\label{sec:helix-content}

This observable measure the percentage of residues that are in a
helical conformation. The boundaries of what is considered a helix can
be adjusted using the four angle settings. The \texttt{per-residue}
option allows you to output the results in a vector format, which can
be read by the \texttt{pdb-b-factor} output mode described above.

\optiontitle{Settings}
\begin{optiontable}
\option{min-angle1}{real}{-1.57}{Minimum angle1 boundary.}
\option{max-angle1}{real}{-0.52}{Maximum angle1 boundary.}
\option{min-angle2}{real}{-1.34}{Minimum angle2 boundary.}
\option{max-angle2}{real}{-0.30}{Maximum angle2 boundary.}
\option{per-residue}{bool}{false}{Output for each residue individually.}
\end{optiontable}


\section{PDB (\texttt{pdb})}
\label{sec:observable-pdb}

Observables can also be used to extract complete structural states
from a simulation. One examples is the PDB observable, which will dump
pdb files are regular intervals. If the output filename contains
special characters \texttt{\%e} (energy value) and/or \texttt{\%i}
(iteration number), the structures will be dumped to individual
files. If not, they will be dumped as models in a single pdb file.

\optiontitle{Settings}
\begin{optiontable}
\option{minimum-energy-mode}{bool}{false}{Whether only to dump minimum energy structures.}
\option{minimum-energy-fraction}{real}{0.1}{How far from the minimum energy a structure must be before being dumped.}
\option{minimum-energy-interval}{int}{500}{Minimum interval between dumped minimum energy structures.}
\end{optiontable}



\section{GIT (\texttt{git})}
\label{sec:observable-git}

It is also possible to dump structural states in GIT vectors. As for
the PDB observable, the vectors can either be dumped to a single file
or to individual files, depending on whether the \texttt{\%e} or
\texttt{\%i} special characters are included in the output file name.





\section{Angle histograms (\texttt{angle-histogram})}
\label{sec:angle-histogram}

The angle histogram observable maintains statistics on the
distribution of angles in the protein. Several of the general
observable settings are important when using this observable: The
statistics will be updated at intervals given by the
\texttt{register-interval} option and dumped at intervals given by the
\texttt{output-interval} option. Optionally, a
\texttt{register-burnin} can be set to skip a number of iterations at
the beginning.



\section{Selecting observables in \texttt{phaistos.cpp}}

In order to use an energy term as an observable, use the
\texttt{observable} keyword, rather than the \texttt{energy} that is
normally used for specifying energies. Just as for moves and energies,
observables can be specified individually

\begin{verbatim}
$ ./phaistos --observable-helix-content --observable-rg
\end{verbatim}

\noindent
or by using a shorthand

\begin{verbatim}
$ ./phaistos --observable helix-content rg
\end{verbatim}

\noindent
Any specific settings can be specified using square brackets as
described in section \ref{sec:command-line-options}. Here is a simple
example using the angle histogram observable:

\begin{verbatim}
$ ./phaistos --observable \
  angle-histogram[output-target:output_%i.txt,\
  register-interval:1,output-interval:10]
\end{verbatim}

\subsection{Special observables: \texttt{@energy-sum} and \texttt{@energy-terms}}
\label{sec:energy-sum-energy}

It is natural to have both the total energy sum and the individual
energy terms as part of your observable. This is supported through two
special observable options: \texttt{@energy-sum} and
\texttt{@energy-terms}. These observables work by looking up the
previously calculated values in the main energy term -- no
recalculation of energies is done.

An additional benefit of this approach is that observable settings can
be provided for the whole group of energies in one go:

\begin{verbatim}
$ ./phaistos --energy profasi \
  --observable @energy-terms[register-interval:1]
\end{verbatim}



\section{\texttt{evaluate\_observable.cpp}}
\label{sec:evaluate-observable}

The \texttt{evaluate\_observable} program is located in the
\texttt{applications} module. It makes it possible to easily evaluate
energies/observables for a given set of PDB files. It follows the
syntax of the \texttt{phaistos.cpp} command line. To evaluate for
instance the helix content of the structure in a PDB file, one would
write:

\begin{verbatim}
$ ./evaluate_observable --observable helix-content \
   --pdb-file 1enh.pdb
# ID    helix-content   
1enh.pdb  0.71153846153846156
\end{verbatim}

\noindent
One distinction from the \texttt{phaistos.cpp} executable is that
\texttt{evaluate\_ob\-serv\-ab\-le} can take many PDB-files at once. This is
done using the \texttt{pdb-files} option:

\begin{verbatim}
$ ./evaluate_observable --observable helix-content \
  --pdb-files *.pdb
# ID    helix-content   
sample1.pdb 0.80769230769230771
sample2.pdb 0.78846153846153844
...
\end{verbatim}

\noindent
Of course, the whole observable machinery is available from
\texttt{evaluate\_ob\-serv\-ab\-le.cpp}, making it possible to determine
where and how often results are outputted. Using the
\texttt{pdb-header} and \texttt{pdb-b-factor} output targets described
above, it is also possible to use \texttt{evaluate\_ob\-serv\-ab\-le}
to annotate a given set of PDB files with energy information.

The \texttt{evaluate\_observable.cpp} program also has an
\texttt{energy} option, which works the same way as for
\texttt{phaistos.cpp}. The output of energies is controlled by the
\texttt{@energy-sum} and \texttt{@energy-terms} observables introduced
in the last section. The default setting for \texttt{--observable} is
to include both \texttt{@energy-sum} and \texttt{@energy-terms}, which
means that if you specify energies and no observables, the sum and
individual terms will be included:

\begin{verbatim}
$ ./evaluate_observable --pdb-file 1enh.pdb --energy opls 
# ID    @energy-sum     clash-fast      opls-angle-bend ...
/PATH/1enh.pdb -3627.4051175360833 0 139.3105497686366 ...
\end{verbatim}

\noindent
Note that if observables are specified, the default will be
overridden, and the energy evaluations might become invisible

\begin{verbatim}
$ ... --energy opls --observable rg                       
# ID    rg      
/PATH/1enh.pdb  10.088942787694526
\end{verbatim}

\noindent
in which case you need to specify the special observables explicitly

\begin{verbatim}
$ ... --energy opls --observable @energy-sum rg
# ID    @energy-sum     rg      
/PATH/1enh.pdb  -3627.4051175360833 10.088942787694526
\end{verbatim}

\chapter[Simulation and Optimization]{Monte Carlo: Simulation and Optimization}
\label{sec:simulation-types}

Phaistos is a Monte Carlo simulation framework. In its broadest sense,
the term Monte Carlo describes any stochastic (non-deterministic)
algorithm. In the context of protein simulations, it is used both to
describe \emph{simulation}, where one wishes to sample from a given
probability distribution (typically the Boltzmann distribution), and
stochastic \emph{optimization}, where the task is merely to find the
structure with lowest energy (highest probability). Simulations are
typically based on a specific MC technique called \emph{Markov Chain
  Monte Carlo} (MCMC), where each new state is generated using only
information about the previous state, and where the property of
\emph{detailed balance} ensures that the simulation will converge to a
given target distribution. In some fields, the Markov Chain Monte Carlo
algorithm is referred to simply by the term Monte Carlo.


\section{MCMC simulations: detailed balance and ergodicity}
\label{sec:detailed-balance}

Markov Chain Monte Carlo is a technique that makes it possible to draw
samples from a complex distribution $\pi$ (target distribution), for
which only the density function is known. The idea is to construct a
simulation where each new state is dependent only on the previous
state. If new states are accepted/rejected in accordance with the
\emph{detailed balance} equation, it can be shown that this Markov
Chain will converge to the target distribution of interest.

The \emph{detailed balance} property entails that the probability of
being in a state $x$ multiplied by the probability of going from state
$x$ to state $x'$, is equal to the probability of being in state $x'$
and moving to $x$

\begin{equation}
\pi(x)P(x \rightarrow x') = \pi(x')P(x' \rightarrow x) \ .
\label{eq:monte-carlo-1}
\end{equation}

\noindent
It can be convenient to factor the transition probability $P(x
\rightarrow x')$ into a \emph{proposal} probability $P_p$ and an
\emph{acceptance} probability $P_a$, that is

\begin{equation}
P(x \rightarrow x') = P_p(x \rightarrow x')P_a(x \rightarrow x') \ .
\label{eq:monte-carlo-2}
\end{equation}

\noindent
In many cases, the moves used in a simulation are symmetric, $P_p(x
\rightarrow x') = P_p(x' \rightarrow x)$. When this is not the case,
the move is said to be \emph{biased}. In such case, the acceptance
probability factors must be chosen in a way to compensate for this
bias. The Metropolis-Hastings method described below corresponds to
one choice of these acceptance probabilities.

It should be noted that the detailed balance equation alone is not
sufficient to guarantee convergence. The set of moves used in the
simulation must also be chosen to ensure \emph{ergodicity}, which
means that there is a non-zero probability of moving from any state to
any other state in a system.


\section{General Monte Carlo settings}
\label{sec:gener-monte-carlo-sett}

In the next sections, the different available Monte Carlo methods in
Phaistos will be described in detail, along with a list of any
available settings. A number of common options are shared between all
Monte Carlo types. To avoid repetition, these are listed here.

\begin{optiontable}
\option{debug}{int}{0}{Debug level.}
\option{declash-on-reinitialize}{bool}{true}{Whether to remove self-collisions in the chain when reinitializing the structure.}
\option{maximum-declash-attempt}{int}{20000}{The maximum number of moves to allow before aborting the declash procedure.}
\option{reinitialization-interval}{int}{0}{How often to reinitialize the simulation (default: 0 (never)).}
\option{consistency-check-interval}{int}{10000}{How often to check whether the simulated system is internally consistent.}
\end{optiontable}



\section{Metropolis-Hastings (\texttt{metropolis-hastings})}
\label{sec:metropolis-hastings}

From expression (\ref{eq:monte-carlo-1}) and (\ref{eq:monte-carlo-2}), we can
write the detailed balance requirement as

\begin{equation}
\frac{P_a(x \rightarrow x')}{P_a(x' \rightarrow x)} = \frac{\pi(x')P_p(x' \rightarrow x)}{\pi(x)P_p(x \rightarrow x')} \ .
\label{eq:monte-carlo-3}
\end{equation}

\noindent
The Metropolis-Hastings method \cite{metropolis1953equation} specifies
that a move from state $x$ to $x'$ should be accepted with the
probability

\begin{equation}
P_a(x \rightarrow x') = \min\left(1, \frac{\pi(x')P_p(x' \rightarrow x)}{\pi(x)P_p(x \rightarrow x')}\right) \ .
\label{eq:monte-carlo-4}
\end{equation}

\noindent
In other words, assuming an unbiased move set, the move is always
accepted if the probability according to the target distribution is
increased ($\pi(x')>\pi(x)$). If the probability decreases, the acceptance
probability depends on how much worse the probability of the new
structure is compared to the previous one. It can be seen that
(\ref{eq:monte-carlo-4}) is a solution to
(\ref{eq:monte-carlo-3}). Other solutions exist, leading to different
convergence properties of the algorithm.  However, the
Metropolis-Hastings criterion is the most frequently used in the
literature.

\section{Greedy Optimization (\texttt{greedy-optimization})}
\label{sec:greedy-optimization}

The simplest Monte Carlo type in Phaistos is greedy optimization. As
in the case for Metropolis-Hastings, a sequence of new states are
generated. However, in this case, only structures improving the energy
are accepted. This can be useful for quickly relaxing a structure.

\section{Simulated Annealing (\texttt{simulated-annealing)}}
\label{sec:simulated-annealing}

The simulated annealing method \cite{kirkpatrick1983optimization} lies
somewhere between Metropolis-Hastings and greedy optimization. It uses
a Metropolis-Hastings acceptance criterion, but rather than using a
fixed temperature, the simulation will gradually move from a high to a
low temperature. This has the effect that initially, the simulation
will accept nearly all generated candidate structures, both with
higher and lower energies. As the temperature decreases, the
simulation will increasingly appear as a greedy optimization,
accepting fewer structures that would increase the energy.

\optiontitle{Settings}
\begin{optiontable}
\option{e-min}{real}{-inf}{Lower bound on energy.}
\option{e-max}{real}{inf}{Upper bound on energy.}
\option{temperature-start}{real}{inf}{Starting temperature (must be finite).}
\option{temperature-end}{real}{inf}{Ending temperature (must be finite).}
\end{optiontable}


% Include module input
\input{/mnt/c/Users/juhasm/Documents/phaistos/modules/muninn/doc/monte_carlo.tex}







\section{Selecting Monte Carlo type in \texttt{phaistos.cpp}}

The type of Monte Carlo algorithm used by \texttt{phaistos.cpp} is set
using the \texttt{--monte-carlo} option. As for energies and moves,
there is a convenient command line shorthand to quickly choose both
the type of Monte Carlo run and the detailed settings:

\begin{verbatim}
$ ./phaistos --monte-carlo metropolis-hastings[debug:1]
\end{verbatim}

\noindent
If no Monte Carlo type is specified, \texttt{phaistos.cpp} will use Metropolis Hastings.



\chapter{Probabilistic Models}
\label{cha:probabilistic-models}

Phaistos contains implementations of a number of probabilistic models,
each modeling distinct aspects of protein geometry. In this chapter,
we will focus on the generative models, which are used both for
proposing new structures, and as energy functions. Other models, which
are used solely as energy functions, are described in separate
appendices.

The generative models are divided into backbone models and sidechain
models. For historical reasons, the backbone models are built into the
core of Phaistos, while the sidechain models are available as a module.

\section{Backbone models}
\label{sec:backbone-models}

The backbone models control the angular degrees of freedom of the
protein main chain. There are several variants: FB5HMM (\texttt{fb5}),
TorusDBN (\texttt{torus}), and TorusCsDBN (\texttt{torus-cs}). The
FB5HMM model is concerned with the $\mathrm{C}_\alpha$ trace of a
protein, while the TorusDB describes the backbone in atomic detail
(N-CA-C, CB, O). Both the FB5HMM and TorusDB models have been
described in detail in the literature \cite{hamelryck2006srp,
  boomsma2008gpm}. The last model, TorusCsDBN, is an extension of the
TorusDBN that includes chemical shift information. In cases where
chemical shift data are available for a system, this can significantly
improve speed and accuracy of a simulation.

\subsection{Using backbone models in \texttt{phaistos.cpp}}
By default, \texttt{phaistos.cpp} will choose a model depending on
your chain type: FB5HMM for $\mathrm{C}_\alpha$ chains and TorusDBN
for full atom chains. You can override this behaviour, or change the
settings of your model using the backbone-dbn option:

\begin{verbatim}
$ ./phaistos --backbone-dbn torus-cs
\end{verbatim}

\noindent
The different models have different associated settings, depending on
the available node types. However, all models have a
\texttt{initial-aa-file} option, an \texttt{initial-ss-file} option,
and an \texttt{initial-pdb-file} option, which allows you to
initialize your model from an amino acid sequence file, a secondary
structure file or a PDB file. In the following command
\texttt{torus-cs} is initialized from the PDB file \texttt{file.pdb}:

\begin{verbatim}
$ ./phaistos --backbone-dbn torus-cs[initial-pdb-file:file.pdb]
\end{verbatim}

\noindent
If no input settings are explicitly specified for the model,
\texttt{phaistos.cpp} will use the global \texttt{pdb-file},
\texttt{aa-file} and \texttt{ss-file} options. The last example
is therefore equivalent to:

\begin{verbatim}
$ ./phaistos --pdb-file file.pdb --backbone-dbn torus-cs
\end{verbatim}

\noindent
For a complete list of all available options, please see the Phaistos
help output:

\begin{verbatim}
$ ./phaistos --help
\end{verbatim}


\section{Sidechain models}
\label{sec:sidechain-models}

As for the backbone case, the two available sidechain models
correspond to two different representations of protein
sidechains. BASILISK is a probabilistic model that represents
sidechains in full atomic detail and in continuous space, modeling all
the individual $\chi$ angles in each sidechain
\cite{harder2010beyond}.  COMPAS uses a more coarse-grained
representation, where sidechains are represented by a single
pseudo-atom. These models are described in greater detail in appendix
\ref{cha:sidechain-dbns-module}. The models are set up through the
energies and moves that use them.




\appendix

\chapter{SAXS module (\texttt{saxs})} \label{cha:saxs-module}

Small angle X-ray scattering (SAXS) is an experimental technique which
can readily provide low resolution structural information on proteins
in solution. The experiment requires relatively low amount of sample:
typical quantities are a volume of $\sim15\,\mu l$, and a sample
concentration of $\sim$1.0\, mg/ml~\cite{SvergFT-1987}.

SAXS was originally developed for structure analysis of condensed
matter in isotropic environments and the first applications date back
to the 1930s, with the publication of studies on metal
alloys~\cite{Guini-1939}. The method was later extended to other
materials, included biological macromolecules in solution~\cite{KratkG-1982}.

The information that can be obtained directly from a SAXS experiment
includes the radius of gyration, volume, molecular mass (provided a
known reference is available) and possibly a coarse-grained
description of the electron distribution in the sample. These
parameters are inferred from inverse Fourier analysis of the SAXS
experiment, and are limited by the noise level of the
instrumentation. In particular, the determination of the electronic
distribution function $p\left(r\right)$ requires a number of
regularization steps whose relative weight is unclear. Manual tuning
is therefore typically required to produce satisfactory
results~\cite{SemenS-1991}.

Since the 1980s, improvement in the instrumentation has led to
reliable analyses of polymeric systems, and increase in computation
power during the 1990s allowed for Bayesian regularization of the
inverse Fourier analysis, providing better estimators and more
rigorous error analysis~\cite{Hanse-1990}.  The latter is
unfortunately still not routinely applied.

Historically, analysis of obtained SAXS curves was limited to a
comparison to scattering curves describing simple geometrical bodies
with known analytical solutions. As sufficient computational power
became available, this analysis was extended with the possibility of
evaluating forward models of SAXS data from a hypothetical
experiment. These structural proposals are typically compared to the
experimental result in order to validate one or more structural
models~\cite{KochVS-2003,KratkG-1982,SvergFT-1987}.

There are two main approaches for calculating a theoretical SAXS curve
from the structure (the forward model). One is based the approximation
of the electronic inhomogeneity via spherical harmonics
expansion~\cite{svergun1995crysol}, the other evaluates the pairwise
interactions occurring between elementary scatterers by means of
direct Fourier transforms~\cite{glatter1982small}.  The spherical
harmonics approach is computationally attractive, but it is also quite
sensitive to the correct determination of the atomic volume of each
individual scatterer. Further, the algorithm enforces a compact,
quasi-globular shape, and the evaluation has to be completely
recomputed even in case of a minor update in the proposal. For this
reason, Monte Carlo protocols do not fully benefit of this method.

The SAXS module in Phaistos follows the description using
pairwise interactions. This is based on the Debye formula~\cite{Debye-1915}

\begin{equation}
I(q)=\sum_{i=1}^{M}\sum_{j=1}^{M}F_{i}(q)F_{j}(q)\frac{\sin(q\cdot r_{ij})}{q\cdot r_{ij}} \ ,
\end{equation}

\noindent where the experimental SAXS curve $I(q)$ is defined as the
sum of the pairwise interactions between the $M$ bodies in the
system. These are individually described by the product of the form
factors, $F_{i}(q)$, with the structure factors depending on the
distance between them ($r_{ij}$) and the scattering momentum
$q=\frac{4\pi}{\lambda}\sin\theta$. Here, $\theta / 2$ is the
scattering angle between the incoming beams, and $\lambda$ the X-ray
wavelength. This approach is computationally more expensive than the
spherical harmonics evaluation ($O(n^2)$) but allows for partial
updates in the structure and hence caching. Importantly, it is not
limited by assumptions on the protein shape.

An attractive performance can be achieved by reducing the number of
elementary scatterers. By grouping neighboring atoms in a set of
\emph{dummy bodies}, the $I(q)$ profile can be conveniently
approximated at least to the current experimental
resolutions~\cite{stovgaard2010calculation}. The Phaistos SAXS module
includes parameters for two levels of coarse-graining. The one-body
model describes each protein residue a single scattering body, whereas
the two-body model used one (common) dummy body for the backbone part
of the residue and one for the side chain scattering. The two-body
model is significantly more accurate for resolutions above $q \approx
0.2 \mathrm{\AA}^{-1}$ while the one-body model evaluates around three
times faster.  A detailed description of the coarse-grained models can
be found in the publication~\cite{stovgaard2010calculation}.

The SAXS module is implemented as an energy term for Phaistos and the
configuration is described in Section~\ref{sec:saxs-energy}.




\bibliographystyle{unsrt}
\bibliography{references.bib\phaistosbibtexfiles}

\end{document}












% \begin{optiondescription}
%   \optionitem[1][int]{move-length-min} Minimum length of move
%   \optionitem[1][int]{move-length-max} Maximum length of move
%   \optionitem[true][bool]{implicit-rotamer-energy} If
%   implicit-rotamer-energy is set to false, the bias introduced when
%   proposing from the Dunbrack library is compensated for, ensuring
%   that the uniform distribution is reached when simulating with no
%   force-field.
%   \optionitem[1.0][real]{sigma-scale-factor} The width of the Dunbrack
%   Gaussian distributions from which new angles are proposed can be
%   scaled by this factor.  Flattening the distribution can be useful to
%   ensure ergodicity, as the Dunbrack library does not have support
%   over the entire interval $[0,2\pi]$
%   \optionitem[1.0][real]{rotamer-state-resample-frequency} Frequency
%   with which move will resample new rotamer states ($g^+, t, g^-$).
%   \optionitem[true][bool]{rotamer-skip-proline} Whether to sample
%   sidechain angles in proline. N.B. If rotamer-skip-proline = false a
%   change in proline bond length is introduced, this must be taken in
%   account by an appropriate bond-stretch energy term.
%   \optionitem[true][bool]{sample-hydrogen-chis} Whether to sample
%   (uniformly) non-standard torsion angles in sidechain (e.g methyl
%   groups, CB-S-H in Cys, etc...)
% \end{optiondescription}

% \newpage

% \begin{small}
% \noindent
% % \begin{tabularx}{\textwidth}{p{0.4\textwidth}llp{0.4\textwidth}}
% \begin{tabular*}{\textwidth}{@{\extracolsep{\fill}}p{0.4\textwidth}llp{0.4\textwidth}}
% % \begin{tabularx}{p{.25\textwidth}llp{0.25\textwidth}}
%   name & type & default & description \\
%   \hline \\
% \end{tabular*}
% \end{small}

% @item -{-}implicit-rotamer-energy
% Boolean (default:true). If implicit-rotamer-energy is set to false, the bias introduced when proposing from the Dunbrack library is compensated,
% therefore ensuring that the uniform distribution is reached when simulating with no force-field.

% @item -{-}sigma-scale-factor
% Float (default:1). The width of the Dunbrack Gaussian distributions from which new angles are proposed can be scaled by this factor. 
% Flattening the distribution can be useful to ensure ergodicity, as th Dunbrack library does not have support over the entire interval $[0,2\pi]$
% @item -{-}rotamer-state-resample-frequency
% Float (default:1, range:0-1). Frequency with which move will resample new rotamer states ($g^+, t, g^-$). 
% @item -{-}rotamer-skip-proline
% Boolean (default: true). Whether to sample sidechain angles in proline. N.B. If rotamer-skip-proline = false a change in proline bond length 
% is introduced, this must be taken in account by an appropriate bond-stretch energy term.
% @item -{-}sample-hydrogen-chis
% Boolean (default:true). Whether to sample (uniformly) non-standard torsion angles in sidechain (e.g methyl groups, CB-S-H in Cys, etc...)




% \tt{--move-lenth} & real & 1.0 & Relative weight of the move \\
% &&&(probability of choosing this move) \\
% \tt{--debug} & int & 0 & Debug level \\
% \tt{--regions} & vector$<$pair$<$int$>>$ & [[0,size]] & Regions affected by move \\
% \end{tabular}





% \item[\tt{--iterations} \emph{n} ] 
%   For how many iterations to run the simulation.

% \newcommand{\optionitem}[3]{\item[{\tt #1} {\tt #2}]  #3}
% \newcommandtwoopt{\optionitem}[3][][]{\item[{\tt #3} {\rm type:#1 default:#2}]}
% \newcommandtwoopt{\optionitem}[3][][]{\item[{\tt --#3 #1} {\rm (=#2)}]}



% \setdescription{labelsep=*}

% \begin{description}
% % \item[--verbose] produce a verbose output (highly recommended if
% %   you want to see what is going on)
% % \optionitem[bool][false]{verbose} hej 
% % \optionitem[bool][false]{verbose} hej 
% % \item[bla] hej
% % \item[bla] hej
% % produce a verbose output (highly recommended if
% %   you want to see what is going on)
% \end{description}


% \begin{description}
% \item[\tt{--weight}] Relative weight of the move (probability of choosing this move)
% \item[\tt{--debug}] print debugging information
% \item[\tt{--move-length-min} ] 
% \item[\tt{--move-length-max}] 
% \end{description}


% \begin{optiondescription}
% \optionitem[1.0]{real}{weight} Relative weight of the move (probability of choosing this move)
% \optionitem[0]{int}{debug} Debug level 
% \optionitem[{[[0,size]]}]{vector$<$pair$<$int$>>$}{regions} Regions in the molecule that will be affected by the move \\
% \end{optiondescription}


% Environment used for option descriptions
% \newenvironment{optiondescription}{
%   \setlength{\leftmargini}{0em}
%   \description
%   \setlength{\itemindent}{0em}
%   \setlength{\labelindent}{0em}
%   \setlength{\labelsep}{\leftmargin-\labelwidth}
% }
% {\endlist}
% \newcommandtwoopt{\optionitem}[3][][]{\item[{\tt-{-}#3 \rm #2} {\rm (=#1)}]\hfill\\}

% \newcommand{\optiontable}{
% \small
% \noindent
% \begin{small}
% \begin{tabularx}{\linewidth}{XllX}
%   name & type & default & description \\
%   \hline \\
% }

  
% \newenvironment{optiontable}{\par\noindent\par\begin{small}\par\begin{tabularx}{\linewidth}{XllX}\par}{\par\end{tabularx}}
% \end{small}
% }
% \newenvironment{optiontable}{\par\begin{tabular}{llll}}{\end{tabular}}
% \newenvironment{optiontable}{\par\begin{tabular*}{\textwidth}{llll}}{\end{tabular*}}
% \end{small}
% }

% \option{--move-length-min} & int & 1
% \newcommand{\optionitem}[3][]{\item[{\tt-{-}#3 \rm #2} {\rm (=#1)}]\hfill\\}



% Item used for option descriptions: \optiondescription[
% \newcommand{\optionitem}[3][]{\item[{\tt-{-}#3 \rm #1} {\rm (=#2)}]\hfill\\}

